% Generated draft: requires mathematical and source review.
\documentclass[11pt]{article}
\usepackage[margin=25mm]{geometry}
\usepackage{fontspec}
\setmainfont{Latin Modern Roman}
\usepackage{amsmath,amssymb,mathtools,unicode-math}
\setmathfont{Latin Modern Math}
\usepackage{enumitem,needspace,xurl,hyperref}
\setlength{\parindent}{0pt}
\setlength{\parskip}{5pt}
\setlength{\emergencystretch}{3em}
\newcommand{\N}{\mathbb{N}}
\newcommand{\Z}{\mathbb{Z}}
\newcommand{\Q}{\mathbb{Q}}
\newcommand{\R}{\mathbb{R}}
\newcommand{\C}{\mathbb{C}}
\newcommand{\F}{\mathbb{F}}
\newcommand{\A}{\mathbb{A}}
\newcommand{\PP}{\mathbb{P}}
\newcommand{\E}{\mathbb{E}}
\newcommand{\eps}{\varepsilon}
\newcommand{\dd}{\,\mathrm{d}}
\begin{document}
\section*{UnsolvedMath ID: 800775}
\textbf{Source problem code:} catalogue entry
\subsection{Definitions and mathematical statement}
For notation, let \(k\) be an algebraically closed field and let
\[
\mathbb F_n=\mathbb P_{\mathbb P^1_k}\bigl(\mathcal O_{\mathbb P^1_k}\oplus\mathcal O_{\mathbb P^1_k}(n)\bigr),\qquad n\in\mathbb Z_{\geq 0},
\]
be a Hirzebruch surface. The supplied record does not specify the ground field or characteristic, or explicitly state the intended range of \(n\). This notation therefore does not supply a missing universal quantifier over fields or indices; the intended setting requires confirmation.

For four pairwise distinct points \(p_1,\ldots,p_4\in\mathbb F_n(k)\), write
\[
\pi:X=\operatorname{Bl}_{\{p_1,\ldots,p_4\}}(\mathbb F_n)\longrightarrow\mathbb F_n
\]
for their blow-up, and let \(T_X\) denote the tangent bundle of \(X\). Under the numerical tautological interpretation, a vector bundle \(E\) on a smooth projective variety is called pseudo-effective if
\[
[\mathcal O_{\mathbb P(E)}(1)]\in\overline{\operatorname{Eff}}^{\,1}(\mathbb P(E)),
\]
where \(\mathbb P(E)\) parametrizes one-dimensional quotients, \(\mathcal O_{\mathbb P(E)}(1)\) is the tautological quotient line bundle, and \(\overline{\operatorname{Eff}}^{\,1}\) is the closure of the cone generated by effective divisor classes in the real numerical divisor space. The supplied record does not define pseudo-effectivity for vector bundles; this interpretation must be checked against the cited source.

\textbf{Problem (Matsumura, 2022, Problem 10).}
\begin{enumerate}
\item Does the blow-up of a Hirzebruch surface at four general points have pseudo-effective tangent bundle? Under the numerical interpretation above, the question asks whether
\[
[\mathcal O_{\mathbb P(T_X)}(1)]\in\overline{\operatorname{Eff}}^{\,1}(\mathbb P(T_X)).
\]
The source does not specify whether ``general'' means points in a nonempty Zariski open subset of the space of distinct ordered quadruples or very general points, nor whether the answer is sought uniformly in \(n\) or separately for each \(n\).
\item What can be said about pseudo-effectivity of the tangent bundle of a Hirzebruch surface blown up at special points? This is an exploratory question: the supplied record specifies neither the special configurations nor the number of points in this second question. It also does not say whether infinitely near centers are allowed. Consequently, a precise classification target cannot be recovered from this record alone.
\end{enumerate}
\subsection{Short English statement}
Matsumura asks whether blowing up a Hirzebruch surface at four general points gives a surface with pseudo-effective tangent bundle, and what happens for blow-ups at special points. The ground field, meaning of general and special points, and precise vector-bundle positivity convention require confirmation. The supplied registry labels the problem open; this draft does not independently verify that status.
\subsection{Sources}
\begin{itemize}
\item Local UnsolvedMath source registry, permanent ID \(800775\), exact source problem code \texttt{catalogue entry}, title ``Positively Curved Projective Varieties --- Problem 10.'' The supplied record attributes the problem to Matsumura, with proposed year 2022, and labels it \texttt{open}; that label is supplied curation, not an independently verified mathematical assessment.
\item Matsumura, \emph{Open Problems on Structure of Positively Curved Projective Varieties} (2022), Problem 10. Supplied source link: \url{https://arxiv.org/abs/2211.09153}. The provenance reports extraction by \texttt{arxiv\_latex} and access on 2026-10-16. These are supplied provenance details; no independent source visit or verification is claimed.
\item Supplied upstream revision: \texttt{d371ee917ac9647fbb06cd6fc529deec6f3dda82}. The registry gives creation and update timestamps of \texttt{2026-10-16T14:00:00Z}; these timestamps do not establish a literature review.
\end{itemize}
\end{document}
